% !TEX root = ../main.tex
\section{Premise: What Generation Does Not Establish}
\label{sec:premise}

One premise carries the rest of the paper, and it is deliberately modest. A
system's producing a statement is a different thing from that statement being
supported by evidence, applicable to the case at hand, checked, and permitted for
use. Any architect will grant the distinction in every other context: a report that
renders is not a report that reconciles, and a value that populates a field is not
a value that passed validation. The premise says only that generated text is
subject to the same distinction, and that the generation event is not what closes
it.

The argument below needs no position on what a model understands or represents
internally. It needs the distinction stated formally, which is the business of this
section.

\subsection{Notation}
\label{sec:notation}

\begin{center}
\begin{tabular}{@{}ll@{}}
\toprule
$M$ & the language model: a parameterised generative function \\
$S$ & the deployed system: $M$ together with its decoding, retrieval, policy,
      and control layers \\
$x$ & an input---a prompt, a structured request, a conversation state \\
$p, q$ & propositions or claims \\
$c$ & the resolved context of a decision \\
$E$ & an evidentiary basis \\
$V$ & a verification procedure \\
$I$ & an institution or accountable organisation \\
\bottomrule
\end{tabular}
\end{center}

The distinction between $M$ and $S$ is load-bearing rather than pedantic. Every
property this paper identifies as necessary for reliance is a property of $S$, and
none is a property of $M$. A model has no version of an audit record, an
authorization, or a resolved context; a system can have all three. Where the
literature reports a result about models, I say $M$; where the object of
governance is at issue, I say $S$.

\subsection{The resolved context}
\label{sec:resolved-context}

A great deal of enterprise information is true only relative to a governing frame.
Write
\begin{equation}
  c = \{\, j,\ t_v,\ f,\ u \,\},
  \label{eq:context}
\end{equation}
where

\begin{center}
\begin{tabular}{@{}l>{\raggedright\arraybackslash}p{10.6cm}@{}}
\toprule
$j$ & jurisdiction or governing regime: the legal, regulatory, or policy authority
      whose rules apply \\
$t_v$ & \emph{valid time}: the effective date the assertion is about---the as-of date
      against which the governing rules are to be assessed \\
$f$ & material case facts: patient facts, customer facts, transaction facts,
      product facts, or whatever else is decision-relevant \\
$u$ & intended use: the decision class, the reliance context, what the answer will
      be used to do \\
\bottomrule
\end{tabular}
\end{center}

A second time dimension is required and it does not belong in $c$. Write $t_k$ for
\emph{transaction time}: the state of the evidential base the system actually had
access to---which corpus release, ingested to what date. The distinction is the
ordinary bitemporal one and it is not decoration.

$t_v$ is a property of the \emph{question}: which date's rules govern this decision.
$t_k$ is a property of the \emph{answer's provenance}: what the system could have
known when it answered. Correctness is assessed against $t_v$; reconstruction requires $t_k$; and their
relation is checkable in \emph{both} directions, which is the whole reason two
timestamps are needed:
\begin{align}
  t_k \prec t_v &\;\Longrightarrow\; \text{the base may omit instruments in force at }
    t_v \quad \text{(omission)}, \label{eq:bitemporal-stale}\\
  t_k \succ t_v &\;\Longrightarrow\; \text{the base may contain instruments not in
    force at } t_v \quad \text{(anachronism)}. \label{eq:bitemporal-anachronism}
\end{align}

Omission is the familiar case: a rule amended before $t_v$ but ingested after $t_k$
leaves an assertion ungrounded in a way no inspection of the retrieved evidence will
reveal, because the missing instrument is missing.

Anachronism is the other half and in regulated work it is the more dangerous of the two,
because the answer is confidently wrong rather than incompletely right. A corpus current
to today, answering a question as-of two years ago, will retrieve provisions that had
not yet been made. Applying a rule to a transaction that preceded it is a familiar and
serious error in its own right, and the system will produce a clean citation to a real
instrument while committing it.

Neither failure is prevented by having a recent corpus, and both are prevented by the
same condition. $\Applicable(E, c)$ must be read to require \emph{interval containment}
on the temporal component of $c$: the instrument's own validity period includes $t_v$.
That excludes the not-yet-enacted and the already-superseded together, and it is a
filter on the evidence rather than a property of the corpus. Temporal validity is
therefore part of applicability rather than a dimension of its own, which is why the
seven conditions of Definition~\ref{def:grounded-evidence} do not list it separately. Both directions are also detectable at the aggregate level
without knowing what any amendment said: compare the corpus ingest watermark against
the decision date, per source class, and flag either sign of the difference. Systems that record only one timestamp cannot
perform the comparison, and most record only one.

$\Resolved(c)$ holds when the components of \eqref{eq:context} are explicit rather
than assumed or inferred without validation. $t_k$ is not resolved by the requester;
it is recorded by the system, and it belongs to I9 and I14 rather than to I10. The distinction between explicit and inferred is the whole of
it: a system that guesses the jurisdiction correctly has not resolved the context,
because nothing in the record establishes that the guess was checked, and an
auditor cannot distinguish a correct guess from a lucky one.

The truth predicate is correspondingly indexed. $\True(p \mid c)$ is truth relative
to a resolved context, and the relation that fails is:
\begin{equation}
  \True(p \mid j_1, t_1, f_1, u_1)
  \nentails
  \True(p \mid j_2, t_2, f_2, u_2).
  \label{eq:no-transfer}
\end{equation}

\subsection{The non-entailments}

Write $\Pr_M(\mathrm{utter}(p) \mid x)$ for the probability that $M$, given input
$x$, produces an assertion with content $p$. The premise is that a high value of
this quantity does not establish any of the properties on which reliance depends:

\begin{align}
  \Pr_M(\mathrm{utter}(p) \mid x) \text{ high}
    &\nentails \True(p \mid c),            \label{eq:ne-true}\\
  \Pr_M(\mathrm{utter}(p) \mid x) \text{ high}
    &\nentails \Supports(E, p, c),         \label{eq:ne-supported}\\
  \Pr_M(\mathrm{utter}(p) \mid x) \text{ high}
    &\nentails \Applicable(E, c),          \label{eq:ne-applicable}\\
  \Pr_M(\mathrm{utter}(p) \mid x) \text{ high}
    &\nentails \Verified(V, p, c),         \label{eq:ne-verified}\\
  \Pr_M(\mathrm{utter}(p) \mid x) \text{ high}
    &\nentails \Warranted_I(p \mid c).     \label{eq:ne-warranted}
\end{align}

Compactly, and in the order in which an institution needs them:
\begin{equation}
  \Pr_M(\mathrm{utter}(p) \mid x) \text{ high}
  \nentails \True(p \mid c)
  \nentails \Supports(E, p, c)
  \nentails \Warranted_I(p \mid c).
  \label{eq:ne-chain}
\end{equation}

These are non-entailments, and the distinction matters because the claim is
frequently misread as stronger than it is. They do not assert that the properties
fail to co-occur. A model can and routinely does produce propositions that are
true, well supported, applicable, verifiable, and fit for reliance. The claim is
that the generation event is not what establishes those properties, so the
generation event cannot be what an institution points to when asked to defend
them. Each arrow in \eqref{eq:ne-chain} marks a gap that something other than the
model has to close.

Nor is the claim that the properties are unrelated. High generation probability is
correlated with truth on the distribution the model was fitted to, sometimes
strongly. Correlation on a training distribution is not entailment, and
Sections~\ref{sec:architecture} and~\ref{sec:canonical-form} concern what happens
when the distribution of interest is not that one.

\subsection{What follows}

The premise is weak, and that is its value: it can be granted without conceding
anything contentious, and it is sufficient for everything in this paper.

Two clarifications, because both are easy to misread from the non-entailments
alone. The claim is not that model output is unreliable. Output is often excellent,
and the argument concerns what can be established about an individual assertion
rather than what can be expected of a population of them. And the claim is not that
the properties in \eqref{eq:ne-chain} are rare. In a well-built system most
assertions have all of them. The point is that they are established elsewhere, by
components that can be pointed to.

That is the whole of the constructive programme. If generation does not establish
support, applicability, verification, or authorization, then those properties are
supplied by something in $S$ that is not $M$---and identifying, building, and
measuring that something is what the rest of this paper is about.
